\begin{afterword}
\summaryhead

The post is about the conjunction fallacy: rating ``A and B'' as more likely than ``A'' alone. In one experiment 68\% of subjects rated ``Reagan will provide federal support for unwed mothers and cut federal support to local governments'' as likelier than the first part alone. The author explains the error by representativeness. An added detail can make an outcome seem more typical, so a story grows more plausible as it grows less probable.

Correcting yourself on a direct comparison is only a patch. In a 1982 experiment, professional forecasters who each saw only one statement still rated a Russian invasion of Poland followed by a suspension of diplomatic relations as likelier than the suspension alone. To avoid that, the author says, they would have had to recoil from the word ``and'', add up improbabilities as log probabilities instead of averaging them, and think of all the reasons the plain event could happen. In a dice experiment, subjects should have chosen the shorter sequence. The general remedy is to feel every added detail as a burden. The post ends by telling readers to question each detail of a futurist's story on its own.

\responsehead

The post's central point is correct. A conjunction is never more probable than its parts, and detail that makes a story more convincing makes it less likely. The point is also Tversky and Kahneman's: the 1983 paper the post cites says that such additions ``can only lower probability'', and quotes the same line from \textsc{The Mikado}. The post cites them in its footnotes for the experiments and for the explanation. What it adds is a remedy, and the remedy is where the problems are.

The remedy is untested. The post asks what the forecasters could have done, and answers, with ``It seems to me'', that they should feel every ``and'' as a shock and every detail as a burden. It reports no evidence that people who do this make fewer errors. Its only illustration of the remedy at work is an anecdote in which the author persuades a listener.

Its two rules of thumb hold only in special cases that the post does not state. The first is to add ``absurdities'', log improbabilities, instead of averaging them. That gives the right answer only when the details are independent. In the stories the post warns against, such as its own China scenario, each step makes the next more likely, so adding the steps' separate absurdities overstates the penalty. The second is to prefer the shorter dice sequence. That works only because the longer sequence contains the shorter. Between sequences that do not contain each other, a shorter one can be much less likely.

The post does get one thing right. The forecasting result compares two groups, and the post offers a fix for each: the group that saw ``and'' could have lowered its estimate, and the group that did not could have raised its own by thinking of reasons.

In short: a correct point from Tversky and Kahneman, with an added remedy that is untested, and rules of thumb that hold only in cases the post does not name.
\end{afterword}
